\documentclass[../../math.tex]{subfiles}
\externaldocument{../../math.tex}
\externaldocument{../basics/natural}

\begin{document}

\setcounter{chapter}{13}

\chapter{Basic Group Theory}

There really isn't much group theory here.  I just proved the main results from
my undergraduate group theory class.  Note that because it is more common to
turn the addition in other types into a group than multiplication, the notation
used in the Rocq code is additive and not multiplicative.  Thus, the notation
used here will be additive as well, despite not being my own personal
preference.

\section{Subgroups}

\begin{definition}
    Let $G$ be a group.  Then we call a set $S : G \to \Prop$ a subgroup of $G$
    if:
    \begin{itemize}
        \item $0 \in S$.
        \item For all $a \in S$ and $b \in S$, we have $a + b \in S$.
        \item For all $a \in S$, we have $-a \in S$.
    \end{itemize}
\end{definition}

\begin{theorem}
    Given a group $G$ and a subgroup $S$, $\T(S)$ is a group, where the
    operation, inverses, and identity are defined to be the same as in $G$.
\end{theorem}
\begin{proof}
    The conditions of a subgroup are precisely the ones needed to allow the
    definitions in $G$ to pass through to $\T(S)$.  Then all of the required
    conditions are trivially true because they are true in $G$.
\end{proof}

\section{Order}

\begin{definition}
    Let $G$ be a group and let $a : G$.  We will define the cyclic subgroup
    $\langle a \rangle$ to be the set $\{na \mid n : \Z\}$.
\end{definition}

\begin{theorem}
    For any $G$ and $a$, the cyclic subgroup is in fact a subgroup.
\end{theorem}
\begin{proof}
    We have $0 = 0a \in \langle a \rangle$, so $0 \in \langle a \rangle$.  When
    $ma \in \langle a \rangle$ and $na \in \langle a \rangle$, we have $ma + na
    = (m + n)a \in \langle a \rangle$.  Finally, when $ma \in \langle a
    \rangle$, we have $-ma = (-m)a \in \langle a \rangle$.
\end{proof}

\begin{definition}
    Let $G$ be a group and let $a : G$.  We will define the order of $a$,
    written $|a|$, to be the cardinality of $\langle a \rangle$.
\end{definition}

\begin{theorem} \label{group_order_countable}
    The order of any value in a group is countable.
\end{theorem}
\begin{proof}
    It suffices to prove that $|\langle a \rangle| \leq |\Z|$.  Every element in
    $\langle a \rangle$ can be written in the form $na$ for some integer $n$.
    Define a function $f : \langle a \rangle \to \Z$ that produces such an
    integer $n$.  Then if $f(x) = f(y)$, we have $x = f(x)a = f(y)a = y$, so $f$
    is injective.
\end{proof}

\begin{theorem} \label{group_order_nz}
    The order of a value in a group is not zero.
\end{theorem}
\begin{proof}
    We have $0 \in \langle a \rangle$, proving that the cardinality of $\langle
    a \rangle$ is at least one.
\end{proof}

In the following, let $a$ be a value in a group $G$, and let $K$ be the set
$\{m : \N \mid m \neq 0 \wedge ma = 0\}$.

\begin{lemma} \label{group_order_fin1}
    If a natural number $n$ is the least value of $K$, then $|a| = n$.
\end{lemma}
\begin{proof}
    Define $f : \N_n \to \langle a \rangle$ by $f(m) = ma$.  We will prove that
    $f$ is bijective.  First, let $f(x) = f(y)$ with $x < n$ and $y < n$.  Then
    $xa = ya$.  Without loss of generality, assume that $x \leq y$.  Then by
    Theorem \ref{nat_le_ex}, we have a natural number $c$ with $x + c = y$.
    Then
    \begin{align*}
        xa &= (x + c)a \\
        xa &= xa + ca \\
        0 &= ca.
    \end{align*}
    Now if $c \neq 0$, we would have $c \in K$, so by the minimality of $n$, we
    have $n \leq c$.  However, $x + c < n$ so $c < n$, a contradiction.  Thus,
    $c = 0$, meaning that $y = x + c = x$, so $f$ is injective.

    To prove that $f$ is surjective, let $ma \in \langle a \rangle$.  Because $n
    \in K$, we have $n \neq 0$, so because the integers are a Euclidean domain
    we have integers $q$ and $r$ such that $m = nq + r$ and $0 \leq r < |n| =
    n$.  Then
    \[
        f(r) = ra = 0 + ra = q0 + ra = q(na) + ra = nqa + ra = (nq + r)a = ma.
    \]
\end{proof}

\begin{theorem} \label{group_order_inf}
    $|a| = |\N|$ if and only if $0 \neq na$ for all nonzero natural numbers $n$.
\end{theorem}
\begin{proof}
    First assume that $|a| = |\N|$.  For a contradiction, assume that there does
    exist a nonzero $n$ with $0 = na$.  Then $K$ is nonempty, so it has a least
    element $m$.  By Lemma \ref{group_order_fin1}, this means that $|a| = m$,
    contradicting $|a| = |\N|$.  Thus, there $0 \neq na$ for all natural numbers
    $n$.

    Now assume that $0 \neq na$ for all natural numbers $n$.  We already know
    that $|a| \leq |\N|$ by Theorem \ref{group_order_countable}, so we only need
    to prove that $|\N| \leq |a|$.  Define $f : \N \to \langle a \rangle a$ by
    $f(n) = na$.  We will prove that $f$ is injective.  Let $f(x) = f(y)$.  Then
    $xa = ya$.  Without loss of generality, assume that $x \leq y$.  Then by
    Theorem \ref{nat_le_ex}, we have a natural number $c$ such that $x + c = y$.
    Then
    \begin{align*}
        xa &= (x + c)a \\
        xa &= xa + ca \\
        0 &= ca.
    \end{align*}
    If $c = 0$, we have $x = x + c = y$ as required.  If $c \neq 0$, we also
    have $ca = 0$, contradicting an assumption of the theorem.  Thus, $f$ is
    injective.
\end{proof}

\begin{theorem} \label{group_order_fin}
    A natural number $n$ is the least value of $K$ if and only if $|a| = n$.
\end{theorem}
\begin{proof}
    The forward direction was Lemma \ref{group_order_fin1} so we only need to
    prove the reverse direction.  Assume that $|a| = n$.  We need to prove that
    $n$ is the least value of $K$.  There are three things to prove: that $|a|
    \neq 0$, that $|a|a = 0$, and that $|a|$ is the least value satisfying these
    two conditions.  $|a| \neq 0$ was precisely Theorem \ref{group_order_nz}.

    To prove that $|a|a = 0$, we will have two further cases for when $K$ is
    empty or nonempty.  If $K$ was not empty, it would have a least value $m$.
    By Lemma \ref{group_order_fin1}, this means that $|a| = m$, and since $ma =
    0$, we have $|a|a = 0$ as well.  If $K$ was empty, by Theorem
    \ref{group_order_inf} we would have $|a| = |\N|$, contradicting $|a| = n$.

    To prove that $|a|$ is the least value in $K$, let $y \in K$, and assume for
    a contradiction that $y < |a|$.  Since $y \in K$, $K$ has a least element
    $m$.  By Lemma \ref{group_order_fin1}, we have $|a| = m$.  However, we now
    have $y < |a|$ and $|a| = m \leq y$, which is impossible.  Thus, we must
    have $|a| \leq y$, showing that $|a|$ is the least value in $K$.
\end{proof}

\section{Lagrange's Theorem}

\begin{definition}
    Let $G$ be a group, $H$ a subgroup of $G$, and $a$ a value in $G$.  Then we
    define the left coset by $a$, denoted $a + H$, as the image of the function
    $f(h) = a + h$ under $H$.  We will say that a set in $G$ is a left coset of
    $H$ if there exsits some $a$ such that the set equals $a + H$.
\end{definition}

\begin{theorem}[Lagrange's Theorem]
    Let $G$ be a group and $H$ a subgroup of $G$.  Let $\mathcal C$ be the
    collection of all left cosets in $G$.  Then
    \[
        |H|\, |\mathcal C| = |G|.
    \]
    In particular, $|H|$ divides $|G|$.
\end{theorem}
\begin{proof}
    Let $g : \mathcal C \to G$ be a function that picks out a value $a$ such
    that $A = a + H$.  So for all $A \in \mathcal C$, we have $A = g(A) + H$.
    Then define a function $f : H \to \mathcal C \to G$ by
    \[
        f(h, A) = g(A) + h.
    \]
    We will prove that $f$ is a bijection.

    For injectivity, let $f(h_1, A) = f(h_2, B)$.  So $g(A) + h_1 = g(B) + h_2$.
    We will first prove that $A = B$.  By the definition of $g$, this is
    equivalent to proving that $g(A) + H = g(B) + H$, which we will do by
    antisymmetry.  First, let $x \in g(A) + H$.  Then $x = g(A) + a$ for some $a
    \in H$.  Then
    \[
        x
        = g(A) + a
        = g(A) + h_1 - h_1 + a
        = g(B) + h_2 - h_1 + a,
    \]
    and since $h_2 - h_1 + a \in H$, we have $x \in g(B) + H$.  Thus, $g(A) + H
    \subseteq g(B) + H$.  Now let $x \in g(B) + H$.  Then $x = g(B) + a$ for
    some $a \in H$.  Then
    \[
        x
        = g(B) + a
        = g(B) + h_2 - h_2 + a
        = g(A) + h_1 - h_2 + a,
    \]
    and since $h_1 - h_2 + a \in H$, we have $x \in g(A) + H$.  Thus, $g(B) + H
    \subseteq g(A) + H$.  By antisymmetry, we have $g(A) + H = g(B) + H$, so $A
    = B$.  Now that $A = B$, from $g(A) + h_1 = g(B) + h_2$, we can cancel to
    get $h_1 = h_2$.  Thus, $(h_1, A) = (h_2, B)$, showing that $f$ is
    injective.

    For surjectivity, let $y : G$.  Define $y' = g(y + H)$.  We know that $y + H
    = y' + H$, so $y \in y' + H$.  This means that there exists an $h : H$ such
    that $y = y' + h$.  Then $-y' + y = h$, showing that $-y' + y \in H$.  Then
    \[
        f(-y' + y, y + H) = g(y + H) - y' + y = y' - y' + y = y,
    \]
    showing that $f$ is surjective.
\end{proof}

\section{Quotient Groups}

\begin{definition}
    Let $G$ be a group, and let $N$ be a subgroup.  Then we call $N$ a normal
    subgroup of $G$ if for all $a : G$ and $h \in N$, $a + h - a \in N$ as well.
\end{definition}

Throughout this whole section, let $G$ be a group and $N$ a normal subgroup.

\begin{definition}
    Define a relation $\sim$ on $G$ by saying that $a \sim b$ if $a - b \in N$.
\end{definition}

\begin{instance}
    $\sim$ is an equivalence relation.
\end{instance}
\begin{proof}
    For any $x : G$, $x - x = 0 \in N$, so $\sim$ is reflexive.  If $a - b \in
    N$, then $-(a - b) = b - a \in N$, so $\sim$ is symmetry.  If $a - b \in N$
    and $b - c \in N$, then $(a - b) + (b - c) = a - c \in N$, so $\sim$ is
    transitive.
\end{proof}

\begin{lemma}
    Addition is well-defined in the quotient by $\sim$.
\end{lemma}
\begin{proof}
    Let $a \sim b$ and $c \sim d$.  We must prove that $a + c \sim b + d$.  So
    we have $a - b \in N$ and $c - d \in N$.  Then because $N$ is a normal
    subgroup, we have $a + c - d - a \in N$.  Then, adding $a - b$, we get
    \[
        a + c - d - a + a - b
        = a + c - d - b
        = a + c - (b + d) \in N,
    \]
    so $a + c \sim b + d$.
\end{proof}

\begin{lemma}
    Additive inverses are well-defined in the quotient by $\sim$.
\end{lemma}
\begin{proof}
    Let $a \sim b$.  We must prove that $-a \sim -b$.  So we have $a - b \in N$.
    Because $N$ is a normal subgroup, we have $-b + a - b + b = -b + a = -(-a -
    -b) \in N$.  Thus, $-a - -b \in N$, so $-a \sim -b$.
\end{proof}

\begin{definition}
    Consider the type $G/{\sim}$.  With addition and additive inverses as given
    by the previous lemmas.  $G/{\sim}$ forms a group called the quotient group
    by $N$, denoted $G/N$.  Furthermore, the map $\phi(x) = [x]$ is trivially
    additive, so it is a group homomorphism from $G$ to $G/N$.
\end{definition}

\begin{theorem} \label{to_qgroup_eq}
    For all $a$ and $b$ in $G$, $a - b \in N$ if and only if $\phi(a) =
    \phi(b)$.
\end{theorem}
\begin{proof}
    This is true by the definition of $\sim$.
\end{proof}

\begin{theorem} \label{to_qgroup_zero}
    For all $a : G$, $a \in N$ if and only if $\phi(a) = 0$.
\end{theorem}
\begin{proof}
    We have $\phi(a) = \phi(0)$, so by the previous theorem this is true if and
    only if $a - 0 = a \in N$.
\end{proof}

\begin{theorem} \label{qgroup_f_ex}
    For all groups $H$ and group homomorphisms $f : G \to H$, if $f(x) = 0$ for
    all $x \in N$, then there exists a group homomorphism $g : G/N \to H$ such
    that $f(x) = g(\phi(x))$ for all $x : G$.
\end{theorem}
\begin{proof}
    First, given $a$ and $b$ such that $a \sim b$, we have $b - a \in N$.  So
    \begin{align*}
        f(b - a) &= 0 \\
        f(b) - f(a) &= 0 \\
        f(b) &= f(a).
    \end{align*}
    Thus, $f$ is well-defined under the equation relation $\sim$, so we can
    transfor $f$ into a group homomorphism $g : G/N \to H$ such that $f(x) =
    g(\phi(x))$ for all $x : G$.
\end{proof}

Whenever possible, the previous theorem will be used when working with group
quotients instead of the general theory of quotients developed previously.

\section{The First Isomorphism Theorem}

In this section, let $G$ and $H$ be groups, and let $f$ be a group homomorphism.

\begin{definition}
    Define the kernel of $f$, denoted $\ker f$, as the set $\{a : G \mid f(a) =
    0\}$.
\end{definition}

\begin{theorem}
    The kernel of $f$ is a normal subgroup of $G$.
\end{theorem}
\begin{proof}
    We have $f(0) = 0$, so $0 \in \ker f$.

    When $f(a) = 0$ and $f(b) = 0$, we have $f(a + b) = f(a) + f(b) = 0$, so $a
    + b \in \ker f$.

    When $f(a) = 0$, then $f(-a) = -f(a) = -0 = 0$, so $-a \in \ker f$.

    Let $f(h) = 0$, and let $a : G$.  Then $f(a + h - a) = f(a) + f(h) - f(a) =
    f(a) + 0 - f(a) = 0$, so $a + h - a \in \ker f$.
\end{proof}

\begin{theorem}
    The image of $f$ is a subgroup of $H$.
\end{theorem}
\begin{proof}
    Because $f(0) = 0$, we have $0 \in f(G)$.

    Because $f(x) + f(y) = f(x + y)$, we have $f(x) + f(y) \in f(G)$.

    Because $-f(a) = f(-a)$, we have $-f(a) \in f(G)$.
\end{proof}

\begin{definition}
    Define a homomorphism $g : G/{\ker f} \to f(G)$ given by Theorem
    \ref{qgroup_f_ex} on $f$ itself.  This is well-defined by the definition of
    the kernel.
\end{definition}

\begin{instance}
    $g$ is bijective.
\end{instance}
\begin{proof}
    For injectivity, let $a$ and $b$ be in $G/{\ker f}$ such that $g(a) = g(b)$.
    We have values $x$ and $y$ in $G$ such that $\phi(x) = a$ and $\phi(y) = b$.
    Then $g(\phi(x)) = g(\phi(y))$, and by the definition of $g$, we have $f(x)
    = f(y)$.  Then $f(x - y) = 0$, so $x - y \in N$.  This means that $\phi(x) =
    \phi(y)$, so $a = b$.  Thus, $g$ is injective.

    For surjectivity, let $y$ be in the image of $f$, so that there exists an
    $x$ such that $f(x) = y$.  Then $g(\phi(x)) = f(x) = y$ by the definition of
    $g$.
\end{proof}

\begin{instance}
    $g^{-1}$ is a group homomorphism.
\end{instance}
\begin{proof}
    \[
        g(g^{-1}(a) + g^{-1}(b))
        = g(g^{-1}(a)) + g(g^{-1}(b))
        = a + b
        = g(g^{-1}(a + b)),
    \]
    and since $g$ is injective,
    \[
        g^{-1}(a) + g^{-1}(b)
        = g^{-1}(a + b).
    \]
\end{proof}

\begin{theorem} \label{group_first_iso}
    $G/{\ker f}$ is isomorphic to $f(G)$.
\end{theorem}
\begin{proof}
    $g$ and $g^{-1}$ are group homomorphisms between the two groups that are
    inverses.
\end{proof}

\end{document}
